\documentclass[11pt,a4paper]{article}
\usepackage[T1]{fontenc}
\usepackage{lmodern,amsmath,amssymb,booktabs,microtype}
\usepackage[margin=24mm]{geometry}
\usepackage[colorlinks=true,urlcolor=blue,linkcolor=blue]{hyperref}
\newcommand{\EE}{\mathbb E}
\newcommand{\tr}{\operatorname{tr}}
\newcommand{\NN}{\mathcal N}
\newcommand{\DD}{\mathfrak D}
\newcommand{\PP}{\mathsf P_h}
\newcommand{\W}{\mathsf W}
\newcommand{\FF}{\mathsf F}
\title{Uniform second moments of finite-time memory fields\\A quantitative bridge from Gaussian chaos to energy functions}
\author{Alan Day \and Chengheng\thanks{AI research assistant (OpenAI).}}
\date{Local research note --- 7 November 2026}
\begin{document}\maketitle
\begin{abstract}
We establish uniform annealed second moments for the energy-pinched
finite-time memory field in the common-coupling complex double-scaled
model. The charged-propagator approximation is first extended to fixed
time derivatives and energy insertions. An even Gaussian Dyson expansion
then isolates one potentially growing neutral term. A finite-parameter
Jacobi recurrence compares coupling-Gaussian Wick powers with deterministic
energy polynomials at a quantitative rate. Its entire expansion applies
also to the unbounded finite-size Gaussian Hamiltonian, and removes that
remaining growth after exact energy pinching. Together with the preceding
conserved-return theorem, the result justifies replacing the projected
memory field by its ordinary-propagation term in two-point functions.
The value and convergence of the limiting fluctuation kernel remain to be calculated.
\end{abstract}

\section{Model and the finite-time problem}
Use standard CAR, $\hbar=1$, string time, and the overlapping
complex $p$-body Gaussian model
\begin{align}
 H&=\sum_{I,J}J_{IJ}c_I^\dagger c_J,\quad J=J^\dagger,&
 \EE J_{IJ}J_{KL}&=w_N\delta_{IL}\delta_{JK},\\
 w_N&=4^p\frac{J_s^2}{\lambda_N}\left(\frac{p!}{N^p}\right)^2,&
 \lambda_N&=p^2/N.
\end{align}
All charge sectors use the same sample $J$. Take even $N\to\infty$,
$\lambda_-\le\lambda_N\le\lambda_+$ with $\lambda_->0$, fixed
$J_s>0$, and compact real memory times $|s|\le S$.
The state is the full $\beta=\mu=0$ normalized annealed trace,
\begin{equation}
 \|Y\|_2^2=\sum_m b_m\EE\tr_mY_m^\dagger Y_m,\qquad
 b_m=2^{-N}\binom Nm,\quad \tr_m=\operatorname{Tr}/\binom Nm.
\end{equation}
At charge $m$ put $h=H_m$, $h_-=H_{m-1}$, and
\begin{equation}
 T(X)=\frac1m\sum_i c_i^\dagger Xc_i,\quad
 F_m(s)=T(e^{ish_-}),\quad \alpha=(1-p/m)_+.
\end{equation}
The variance and a uniform bound are
\begin{equation}
 C=C_m=w_N\binom mp\binom{N-m+p}p,\qquad
 M=\sup_m C_m\le\frac{J_s^2}{\lambda_N}e^{2p/N}.
\end{equation}
The covariance map $\Phi(X)=\EE_u[uXu]$ always uses a new copy
$u$ in this charge, with its input held fixed.

\newpage
\section{Memory fields and the result}
For $m>0$, use the exact charge-removal dilation
\begin{align}
 V\psi&=m^{-1/2}\sum_i|i\rangle c_i\psi,&
 K&=1_N\otimes H_{m-1},&Q&=1-VV^\dagger,\\
 a&=V^\dagger KV=\alpha h,&R&=QKV,&D&=QKQ,\\
 M_m(s)&=R^\dagger e^{isD}R,&
 \Lambda_m(s)&=R^\dagger e^{isK}R.
\end{align}
For $m<p$ the force vanishes, so clipping $\alpha$ changes nothing.
Energy pinching retains complete degenerate blocks:
\begin{equation}
 \PP Y=\sum_eP_eYP_e.
\end{equation}
Define the actual and ordinary-propagation fields
\begin{align}
 \FF_m(s)&=(1-\PP)m e^{-ihs/2}M_m(s)e^{-ihs/2},\\
 X_{0,m}(s)&=(1-\PP)m e^{-ihs/2}\Lambda_m(s)e^{-ihs/2}.
\end{align}
The earlier return calculation \cite{Return} establishes
\begin{equation}
 d_N(S):=\sup_{|s|\le S}\|\FF(s)-X_0(s)\|_2\longrightarrow0.
 \label{eq:priorreturn}
\end{equation}
It did not by itself bound either second moment. The new conclusion is
\begin{equation}
 \boxed{\quad
 \sup_N\sup_{|s|\le S}\|X_0(s)\|_2<\infty,\qquad
 \limsup_{N\to\infty}\sup_{|s|\le S}\|\FF(s)\|_2<\infty.
 \quad}
 \label{eq:main}
\end{equation}
The first supremum is along the specified double-scaling sequence;
finitely many initial sizes can be absorbed into its constant.

For $Y=\FF$ or $X_0$, define
\begin{equation}
 \mathcal C_Y(s_1,s_2;t)
 =\EE\tau\big[Y(s_1)^\dagger e^{itH}Y(s_2)e^{-itH}\big].
\end{equation}
Equations \eqref{eq:priorreturn} and \eqref{eq:main} imply
\begin{equation}
 \sup_{|s_1|,|s_2|\le S,\ t\in\mathbb R}
 |\mathcal C_{\FF}(s_1,s_2;t)-\mathcal C_{X_0}(s_1,s_2;t)|\to0.
 \label{eq:cov}
\end{equation}
This is an equivalence of the two finite-$N$ correlators, not an
evaluation or a proof of convergence of their common large-$N$ limit.

The proof has three parts: weighted transfer to an explicit Gaussian
OU force; a finite Dyson expansion retaining its common inserted copy;
and a quantitative, exactly conserved energy reference for the neutral
covariance sandwich. The last part does not assume a finite-size
Hamiltonian has compact spectrum.

\newpage
\section{Time derivatives and energy insertions}
On Gaussian matrix-valued $L^2$ let
$\|\psi\|_r=\|(\NN+1)^{r/2}\psi\|_2$, $r\ge0$. This is a
Gaussian Sobolev norm, not a Schatten-$r$ norm. For
$H=\sum_\beta g_\beta B_\beta$, $\sum B_\beta^2=CI$, creation and
annihilation give, for either left or right multiplication,
\begin{equation}
 \|H_{L/R}\psi\|_r\le A_r\sqrt C\|\psi\|_{r+1},\qquad
 A_r=1+2^{r/2}.
\end{equation}
The number-weight energy estimate gives
\begin{equation}
 \|e^{itH_L}\psi\|_r\le e^{q_r\sqrt C|t|}\|\psi\|_r,
 \quad q_r=r2^{\max(r-1,0)},\quad q_0=0.
 \label{eq:sobprop}
\end{equation}
For example, its adjacent-block estimate is
$\sqrt{k+1}|(k+2)^r-(k+1)^r|/
[(k+1)^{r/2}(k+2)^{r/2}]\le q_r$.

Recall $T=\mathcal U S_-=S_+\mathcal U$ from \cite{Transfer}, and put
\begin{equation}
 n=N-m+1,\quad B_m=\lambda_N+\lambda_N^2+
 \lambda_N(m-n)^2/N,\quad G=\mathcal U H_{-,L}-H_{+,L}\mathcal U.
\end{equation}
In central sectors $N/4\le m\le3N/4$, take $N\ge16$, $p\le N/16$.
The low-chaos defect is bounded by
$10^6\sqrt{CB_m}(k+1)^{3/2}/N$. In the complementary block,
$(k+1)p>N/8$ implies $N\le64\lambda_N(k+1)^2$.
Thus the coarse Gaussian-field bound absorbs the entire high block,
and, with harmless loose constants,
\begin{align}
 \|G\psi\|_r&\le10^7 A_r\frac{\sqrt{MB_m}}N\|\psi\|_{r+5},\\
 \|[\DD,H_L]\psi\|_r
 &\le10^4 A_r\sqrt{M(\lambda_N+\lambda_N^2)}\|\psi\|_{r+3},
 \qquad \DD=p\NN-\mathcal R.
 \label{eq:weightedcomm}
\end{align}
The second line is the all-degree rank estimate \cite{Rank} with
number weights. For $\delta_j(s)=\mathcal Ue^{ish_-}h_-^j-e^{ish}h^j$,
the initial difference is exactly
$\sum_{a<j}H_{+,L}^{j-1-a}G H_{-,L}^aI$.
Weighted Duhamel using \eqref{eq:sobprop} therefore controls
$\delta_j$ by $K_{r,j,S,M}\sqrt{B_m}/N$.
Similarly, starting from
$\DD H^jI=\sum_{a<j}H_L^a[\DD,H_L]H_L^{j-1-a}I$ gives
$\sup_{|s|\le S}\|\DD h^je^{ish}\|_r\le K\sqrt{\lambda_N+\lambda_N^2}$.

Apply the preceding exact rank-filter inequality to $h^je^{ish}$:
\begin{equation}
 \|(S_+-\alpha^{\NN})\phi\|_r\le
 \frac4N\|\DD\phi\|_r+
 \frac{8p|m-n|}{N^2}\|\phi\|_{r+2}+
 \frac{528p^2}{N^2}\|\phi\|_{r+4}.
\end{equation}
For all fixed $r,j,a,b$, the actual derivative estimate is consequently
\begin{equation}
 \boxed{\quad
 \sup_{|s|\le S}\|h^a\partial_s^j[F_m(s)-\mathcal O_\alpha(e^{ish})]h^b\|_r
 \le K_{r,j,a,b,S,M}\frac{\sqrt{B_m}}N.\quad}
 \label{eq:weightedtransfer}
\end{equation}
Here $\mathcal O_\alpha=\alpha^{\NN}$. Finite Hermite polynomials
provide a common core; finite-chaos energy estimates, closure and
Duhamel give the stated domains. No derivative of an unweighted
error estimate was taken.

\newpage
\section{An even Gaussian Dyson expansion of the force}
The exact derivative identity and \eqref{eq:weightedtransfer} give
\begin{align}
 \Lambda&=-F_m''+i\alpha\{h,F_m'\}+\alpha^2hF_mh,\\
 \|\Lambda-\Lambda_{\rm OU}\|_2&\le K_S\sqrt{B_m}/N,\label{eq:forceerror}\\
 \Lambda_{\rm OU}(s)&=\epsilon\EE_u[u e^{is(\alpha h+\sqrt\epsilon u)}u],
 &\epsilon&=1-\alpha^2.
\end{align}
The same $u$ occurs inside and outside the exponential. Fix $h$ and
write $U_t=e^{i\alpha ht}$. The function
$f(t)=\EE_u[u e^{is(\alpha h+tu)}u]$ is even. Taylor expansion in
$t=\sqrt\epsilon$, keeping $\alpha h$ exact, yields
\begin{equation}
 \Lambda_{\rm OU}(s)=\epsilon\Phi(U_s)-\epsilon^2\mathcal I_s(h)
                       +\mathcal E_s,
 \label{eq:dyson}
\end{equation}
where, for either sign of $s$,
\begin{equation}
 \mathcal I_s(h)=s^2\int_{0\le y\le x\le1}
 \EE_u[uU_{s(1-x)}uU_{s(x-y)}uU_{sy}u],dx,dy.
\end{equation}
All three Wick pairings are retained. For unitary $A,B,D$,
\begin{align}
 \EE_u[uAuBuDu]
 ={}&\Phi(A)B\Phi(D)+\Phi(A\Phi(B)D)\notag\\
 &+\sum_{\beta,\gamma}B_\beta A B_\gamma B B_\beta D B_\gamma.
\end{align}
Each term has operator norm at most $C^2$. For the last term use
row--column Cauchy--Schwarz and
$\sum_\beta\Phi(BB_\beta D)^\dagger\Phi(BB_\beta D)\le C^3I$.
It follows that
\begin{equation}
 \|\mathcal I_s(h)\|_\infty\le\tfrac32 C^2s^2.
\end{equation}

For completeness, Gaussian integration by parts and weighted Young's
inequality in a spectral basis of $u$ give
\begin{equation}
 \EE\tr u^{2q}\le(2q-1)C\EE\tr u^{2q-2}
 \le(2q-1)!!C^q.
\end{equation}
This is the scalar-variance case of standard Gaussian matrix moment
control; compare \cite{Tropp}. Every fourth Duhamel derivative of
$f$ contains six factors $u$ and otherwise real-time unitaries.
Normalized Schatten-$12$ H\"older, followed by Jensen, therefore gives
the dimension-free, conditional-on-$h$ bound
\begin{equation}
 \boxed{\quad
 \|\mathcal E_s\|_2\le
 \frac{\epsilon^3|s|^4}{24}\sqrt{11!!}\,C^3.\quad}
 \label{eq:dysonerror}
\end{equation}
In central sectors $m\epsilon\le2p$, $m\epsilon^2\le16\lambda_N$,
and $m\epsilon^3\le128\lambda_N^2/p$. Thus the second term of
\eqref{eq:dyson} is uniformly bounded at collective scale and the
remainder vanishes. Only $p(1-\PP)\Phi(U_s)$ can still grow.

\newpage
\section{A finite-parameter annihilation recurrence}
We now control that neutral term. All statements in this section
are in a central charge sector. Let
\begin{equation}
 Z_j=\W_j(H^j),\qquad
 \mu(x)=\prod_{\ell=0}^{p-1}
 \left(1-\frac{x}{m-\ell}\right)_+
 \left(1-\frac{x}{N-m+p-\ell}\right)_+,
 \quad a_j=\sum_{v=0}^{j-1}\mu(vp).
\end{equation}
$\W_j$ is the coupling-Gaussian chaos projection, not spectral
orthogonalization. The deterministic covariance satisfies
$\Phi(P_rX)=C\mu(r)P_rX$, and $\mu$ is globally
$12p/N$-Lipschitz. With Gaussian creation $\mathcal C$ and
annihilation $\mathcal A=\mathcal C^\dagger$, $Z_j=\mathcal C^jI$.
Differentiating the homogeneous matrix polynomial before projecting gives
the exact identity
\begin{equation}
 \mathcal A Z_j
 =\W_{j-1}\left[\sum_{v=0}^{j-1}\Phi(H^v)H^{j-1-v}\right]
 =Ca_jZ_{j-1}+r_j.
 \label{eq:annih}
\end{equation}
No Wick contractions are omitted; all derivative positions are the
$v$-sum in \eqref{eq:annih}.

Let $\pi_v$ be the positive rank law of a product of $v$ independent
Hamiltonians and $D_v=vp-r$. The earlier all-degree estimate is
\begin{equation}
 \EE_{\pi_v}D_v^2\le2\cdot10^6(\lambda_N v^2+\lambda_N^2v^4).
\end{equation}
For a fixed derivative position, apply the Lipschitz estimate to
the first $v$ independent tensor factors. The remaining factors
are eliminated in the norm using $\EE H_i^2=CI$. Gaussian
symmetrization is a contraction. Hence
\begin{equation}
 \boxed{\quad
 \|r_j\|_2\le\frac{12p}{N}\sqrt{(j-1)!\,C^{j+1}}
             \sum_{v<j}\sqrt{\EE_{\pi_v}D_v^2}.\quad}
 \label{eq:remainder}
\end{equation}
There is only one $(j-1)!$. This estimate is uniform in degree,
including $vp$ beyond the physical rank range.

Retaining finite-$N$ coefficients is useful. Bernoulli's inequality
in each positive factor gives
\begin{equation}
 \mu(jp)\le\mu(p)^j,\qquad
 \mu(p)\le\exp\!\left[-\frac{4\lambda_N}{1+p/N}\right].
\end{equation}
The second estimate uses the reciprocal sum of the two denominators,
which is at least $4/(N+p)$. Thus for all $j\ge1$,
\begin{equation}
 1\le a_j\le A,\qquad A=(1-e^{-3\lambda_-})^{-1},
\end{equation}
under $p\le N/16$. We do not replace $\mu(jp)$ by its leading
$q^j$ limit in the quantitative argument.

\newpage
\section{Auxiliary Jacobi polynomials and the chaos error}
Define monic scalar polynomials
\begin{equation}
 Q_0(x)=1,\quad Q_1(x)=x,\qquad
 Q_{j+1}(x)=xQ_j(x)-Ca_jQ_{j-1}(x).
 \label{eq:Q}
\end{equation}
They belong to an auxiliary positive Jacobi problem \cite{DLMF},
whose off-diagonal entries are $\sqrt{Ca_j}$ and whose diagonal
is zero. It has norm at most $2\sqrt{CA}$.
This auxiliary measure is not asserted to be the actual finite-$N$
annealed density of states.

Put $E_j=Z_j-Q_j(H)$. The leading coupling degree cancels, and
parity leaves degree at most $j-2$. Equations \eqref{eq:annih}
and \eqref{eq:Q} give
\begin{equation}
 E_{j+1}=HE_j-Ca_jE_{j-1}-r_j,
 \qquad \|HE_j\|_2\le2\sqrt{C(j-1)}\|E_j\|_2.
\end{equation}
Write $\rho=p/N$ and
\begin{equation}
 e_j=\frac{\|E_j\|_2}{\sqrt{j!}\,C^{j/2}},\quad
 L=12\sqrt{2\cdot10^6}(\sqrt{\lambda_+}+\lambda_+),\quad K_*=4+A.
\end{equation}
The rank moment bound and \eqref{eq:remainder} imply
\begin{equation}
 e_{j+1}\le2e_j+\frac{A}{\sqrt{j(j+1)}}e_{j-1}+\rho L j^2.
\end{equation}
Since $e_0=e_1=e_2=0$, a direct induction using $j^2\le3^j$
and $K_*\ge5$ yields
\begin{equation}
 \boxed{\quad
 \|Z_j-Q_j(H)\|_2
 \le\rho L K_*^j\sqrt{j!}\,C^{j/2}.\quad}
 \label{eq:bridgepoly}
\end{equation}
Also $\|Z_j\|_2\le\sqrt{j!}\,C^{j/2}$ by independent-tensor
symmetrization. No uniform lower bound on $C$ is required;
the case $C=0$ is simply $H=0$.

Let $J_{\rm aux}$ be the bounded Jacobi operator and define
\begin{equation}
 h_j=C^j\prod_{\ell=1}^j a_\ell,\qquad
 c_j(s)=\frac{\langle e_j,e^{isJ_{\rm aux}}e_0\rangle}{\sqrt{h_j}}.
\end{equation}
Bandedness implies that the numerator's Taylor series starts at
degree $j$. Since $h_j\ge C^j$, for $|s|\le S$,
\begin{equation}
 |c_j(s)|\le e^{2\sqrt{CA}S}\frac{(2\sqrt A S)^j}{j!}.
 \label{eq:coeff}
\end{equation}
This factorial decay beats the $\sqrt{j!}$ in
\eqref{eq:bridgepoly}. It is the required control of the order of
limits: the full degree sum is controlled before the large-$N$ limit.

\newpage
\section{Entire expansion and an exactly conserved reference}
The Jacobi expansion is valid outside its auxiliary spectral interval.
The recurrence implies
$|Q_j(z)|\le(|z|+\sqrt{CA})^j$. Together with \eqref{eq:coeff},
this makes $f_s(z)=\sum_jc_j(s)Q_j(z)$ normally convergent on every
complex compact set. The Jacobi equation gives
\begin{equation}
 c_j'(s)=i[c_{j-1}(s)+Ca_{j+1}c_{j+1}(s)],\qquad c_{-1}=0.
\end{equation}
Termwise differentiation and \eqref{eq:Q} yield
$\partial_s f_s(z)=izf_s(z)$, $f_0(z)=1$. Therefore
\begin{equation}
 e^{isz}=\sum_{j=0}^\infty c_j(s)Q_j(z),\qquad z\in\mathbb C.
\end{equation}
It applies to every finite Gaussian sample's spectrum. Separately,
\eqref{eq:bridgepoly} implies absolute convergence in the actual
annealed matrix $L^2$. Thus no spectral cutoff has been imposed on $H$.

One explicit uniform summation constant is
\begin{equation}
 R_S=\sqrt2 L\exp\!\left(2\sqrt{AM}S+4K_*^2AM S^2\right),
\end{equation}
using $\sum_jx^j/\sqrt{j!}\le\sqrt2e^{x^2}$. Consequently
\begin{equation}
 \left\|e^{isH}-\sum_jc_j(s)Z_j\right\|_2\le\rho R_S.
\end{equation}
For any deterministic sequence $|\theta_j|\le1$, define
\begin{equation}
 g_{\theta,s}(H)=\sum_j\theta_jc_j(s)Q_j(H).
\end{equation}
This is an actual samplewise energy function, not an assertion that
a chaos multiplier commutes with $H$. Contraction of the multiplier
and a second use of \eqref{eq:bridgepoly} give
\begin{equation}
 \boxed{\quad
 \|\theta(\NN)e^{isH}-g_{\theta,s}(H)\|_2\le2\rho R_S.
 \quad}
 \label{eq:multiplier}
\end{equation}
$\theta$ may depend on $N,m,s$, but not on the original random $J$.

Now choose $\theta_j=\mu(jp)$. The earlier rank--chaos filter estimate is
\begin{equation}
 \|(\Phi/C-\theta(\NN))e^{isH}\|_2
 \le12\rho\|\DD e^{isH}\|_2.
\end{equation}
Let $D_S$ denote its already established uniform bound on
$\|\DD e^{isH}\|_2$ \cite{Rank}. Since $g_{\theta,s}(H)$ commutes
with $H$ and pinching is an orthogonal projection, we obtain
\begin{equation}
 \boxed{\quad
 \sup_{|s|\le S}p\|(1-\PP)\Phi(e^{isH})\|_2
 \le C\lambda_N(12D_S+2R_S).\quad}
 \label{eq:neutral}
\end{equation}
This is the needed quantitative energy subtraction. The auxiliary
Jacobi and coupling-Gaussian constructions were never identified
exactly at finite $N$; their difference was explicitly bounded.

\newpage
\section{Assembly of the true memory bound}
Within a central sector, combine \eqref{eq:forceerror},
\eqref{eq:dyson}, \eqref{eq:dysonerror} and \eqref{eq:neutral}.
The phases $e^{-ihs/2}$ are unitary and commute with the pinching.
As $m\epsilon/p\le2$, one obtains a bound of the form
\begin{align}
 \sup_{|s|\le S}\|X_{0,m}(s)\|_2
 \le{}&K_S\sqrt{B_m}
       +2M\lambda_N(12D_S+2R_S)\notag\\
      &+24\lambda_NM^2S^2
       +\frac{128\lambda_N^2}{p}\frac{S^4\sqrt{11!!}M^3}{24}.
 \label{eq:assemble}
\end{align}
The first term is the collective transfer defect, and the third
is the four-$u$ correction. Both can be order one; neither has been
discarded. The last term tends to zero.

The full-charge average satisfies
\begin{equation}
 \sum_m b_m(m-n)^2=N+1,qquad
 \sum_{m\ {
m central}}b_m B_m
 \le\lambda_N+\lambda_N^2+\lambda_N(1+1/N).
\end{equation}
Squaring \eqref{eq:assemble} and summing therefore stays bounded.
For the exterior sectors, $T$ commutes with the Gaussian number
operator and is an $L^2$ contraction. The all-sector variance bound
and \eqref{eq:sobprop} give uniform coarse bounds on the fixed
derivatives and energy insertions in $\Lambda$; multiplying by
$m\le N$ costs at most $N^2$ in the squared norm. The exterior
binomial weight is at most $2e^{-N/8}$, so this contribution vanishes.
This proves the first assertion of \eqref{eq:main}; the second follows
from \eqref{eq:priorreturn}.

Finally Hilbert--Schmidt Cauchy--Schwarz gives, uniformly in real $t$,
\begin{align}
 |\mathcal C_{\FF}-\mathcal C_{X_0}|
 \le{}&\|\FF(s_1)-X_0(s_1)\|_2\|\FF(s_2)\|_2\notag\\
 &+\|X_0(s_1)\|_2\|\FF(s_2)-X_0(s_2)\|_2.
\end{align}
The right side is bounded by a constant times $d_N(S)$, proving
\eqref{eq:cov}.

\paragraph{Next calculation.}
Uniform second moments now hold for the actual finite-memory field.
The remaining task is the joint limit of the bounded components:
the charged-transfer defect, the neutral covariance-filter residual
and the four-$u$ correction, with the exact common disorder and
energy pinching retained. Their cross terms can contribute at order
one. One must compute these terms to determine the finite-time
fluctuation kernel and compare it to the Susskind meson dictionary.
The present bound alone does not supply that kernel or a finite-$G$
bulk interaction law.

\newpage
\section{Verification and sources}
Three fresh scientific outputs accompany the proof. Complete
Gaussian-Hermite coefficients in three CAR systems verify twelve
annihilation recurrences through degrees $5$, $4$ and $3$, including
the lower-chaos terms of $Q_j(H)$. An independent covariance-word
construction checks \eqref{eq:annih}; positive rank tensors check
the single-factorial remainder normalization. A full four-coordinate
Hermitian $2\times2$ Gaussian fixture checks three fourth-moment
Wick words and nine noncommuting finite-time Dyson remainders.
Two Gaussian and time quadrature orders were compared. These
order comparisons are not certified quadrature enclosures.

The auxiliary Jacobi series was also evaluated at $x=\pm6$, outside
its spectral bound, for three finite parameter sets at 80 and 120
digits. Six checks and 105 exact geometric-weight inequalities are
saved. These confirm finite formulas and cancellation conventions;
the uniform bounds and entire continuation are established by the
arguments above, not by finite fits. Inputs, environments, hashes,
build logs and rendered pages are saved locally. Old scientific
outputs are unchanged.

\paragraph{Literature boundary.}
Gaussian matrix moments, Jacobi theory and $q$-Gaussian conditional
expectation are established tools \cite{Tropp,DLMF,BKS}. The
spectral Wick polynomials in \cite{Tall} are defined by annealed
spectral orthogonality. They are not identical at finite $N$ to
$\W_j(H^j)$; \eqref{eq:bridgepoly} supplies the missing comparison
for this argument, using an auxiliary finite-parameter recurrence.
We do not infer priority from a literature search. The new local
result concerns the specified microscopic field and its second
moments, with the prior-stage estimates listed below as dependencies.

\begin{thebibliography}{9}
\bibitem{Tropp} J. A. Tropp, \emph{Second-Order Matrix Concentration
Inequalities}, Prop.~2.3 and Sec.~7,
\href{https://tropp.caltech.edu/papers/Tro16-Second-Order-Matrix-preprint.pdf}{author-hosted preprint} (2016).
\bibitem{DLMF} NIST Digital Library of Mathematical Functions,
\href{https://dlmf.nist.gov/18.2}{Sec.~18.2}, general orthogonal polynomials.
\bibitem{BKS} M. Bo\.{z}ejko, B. K\"ummerer and R. Speicher,
\emph{$q$-Gaussian processes: non-commutative and classical aspects},
\href{https://arxiv.org/abs/funct-an/9604010}{arXiv:funct-an/9604010},
Prop.~2.9 and Theorem~2.11.
\bibitem{Tall} J. Tall and S. Tomsovic, \emph{Double-scaled bosonic
and fermionic embedded ensembles, complex SYK, and the dual Hilbert space},
\href{https://arxiv.org/abs/2604.14522}{arXiv:2604.14522v2}, Sec.~3.1.
\bibitem{Rank} A. Day and Chengheng, \emph{Uniform rank bounds for
finite-time propagation}, local research note (7 November 2026).
\bibitem{Transfer} A. Day and Chengheng, \emph{Adjacent-charge
propagation as Gaussian smoothing}, local research note (7 November 2026).
\bibitem{Return} A. Day and Chengheng, \emph{The conserved background
of the first memory return}, local research note (6 November 2026).
\end{thebibliography}
\end{document}
